\section{Введение}

Инференс больших языковых моделей становится одной из важных вычислительных задач: размер весов, длина контекста и число одновременных запросов предъявляют различные требования к памяти, вычислительным ресурсам и обмену данными. Если модель нельзя эффективно выполнить на одном устройстве, применяют модельный параллелизм, распределяя разные части её вычислительного графа между исполнителями. При послойном разбиении последовательные группы блоков Transformer закрепляются за разными узлами, а промежуточные представления передаются между ними. Такой подход необходимо отличать от параллельной обработки запросов несколькими полными копиями модели \cite{petals2023,deepspeed2022}.

Распределение вычислений само по себе не гарантирует ускорения. Передача активаций, сериализация тензоров, удалённые вызовы, хранение весов и недостаточная загрузка этапов могут перевесить потенциальную выгоду. Исследования последних лет показывают существенную роль коммуникационных механизмов и баланса этапов даже в высокопроизводительных системах на нескольких GPU \cite{jekel2026,hwang2026}. В малых CPU-конфигурациях влияние накладных расходов может быть особенно заметно. При этом мощный GPU-кластер не является необходимым условием для исследования \emph{корректности алгоритма, стоимости информационного обмена и предсказательной модели размещения слоёв}.

В качестве отправной точки используется ранее созданный автором учебно-экспериментальный прототип распределённого инференса \code{sshleifer/tiny-gpt2}, реализованный с помощью PyTorch, Hugging Face Transformers и Ray Actors. В нём модель разделена на три функциональные части и последовательно исполняется в трёх контейнерах Docker Compose на \emph{одной физической машине}. Прототип подтверждён запуском Ray-кластера, проверкой размещения исполнителей, генерацией текста и автоматическими тестами. Однако сама tiny-модель имеет лишь два Transformer-блока и служит проверке архитектуры, а не оценке промышленного многоузлового инференса.

Наиболее конкретное обнаруженное ограничение исходной реализации связано с возвратом полных выходных логитов. При жадном декодировании для продолжения последовательности требуется выбрать один следующий токен по логитам \emph{последней позиции}. Возврат логитов для всех позиций создаёт избыточный объём данных. Но важно установить, какая доля возможной экономии связана с сокращением \emph{передачи}, а какая~--- с уменьшением \emph{вычислений} при формировании логитов. Поэтому простое сравнение «исходная реализация~--- перенос \code{argmax}» недостаточно для причинно интерпретируемого эксперимента.

Другой перспективный аспект~--- выбор границ разделения модели. Для реальных многослойных моделей число возможных разбиений быстро растёт. При неоднородной скорости узлов, разных объёмах доступной памяти и стоимости передачи данных равномерное деление по числу слоёв не обязательно минимизирует задержку. В то же время для одного последовательного запроса нельзя без дополнительных допущений подменять минимизацию \emph{суммарной задержки} минимизацией времени самого медленного этапа; последняя задача непосредственно связана с производительностью перекрывающегося многозапросного конвейера \cite{hwang2026}.

\textbf{Цель исследования}~--- разработать воспроизводимую методику анализа коммуникационных затрат послойного инференса и сформулировать ресурсно-ограниченный способ выбора размещения компонентов модели. В рамках текущего этапа решаются задачи: (1) формализация объёмов обмена и режимов обработки логитов; (2) постановка гипотез о влиянии изменений на задержку; (3) формулировка задачи разбиения с ограничениями памяти; (4) проектирование алгоритма и контрольных проверок; (5) описание авторских нагрузочных данных и реалистичного эксперимента на доступном оборудовании.

\textbf{Границы выполненного и планируемого.} Исходный вариант с полным возвратом логитов уже реализован в НИР. Сокращённые режимы, алгоритм выбора границ разбиения и серия сравнительных экспериментов в настоящей статье предлагаются к реализации. Результаты ускорения, оптимальность на реальном оборудовании и показатели многомашинного кластера не объявляются полученными.
